\section{Engineering Memory Manifesto}

\subsection{Purpose}

This project exists to address a fundamental limitation in modern engineering workflows: the inability to preserve and communicate engineering understanding as projects evolve.

Engineering projects are not defined solely by their source code, documentation, or design artifacts. They are driven by a continuous sequence of decisions, assumptions, experiments, trade-offs, and evolving understanding. Most of this knowledge exists only in the engineer's mind and gradually disappears over time.

The objective of this project is to preserve this engineering understanding in a structured, continuously evolving form.

\vspace{0.5cm}

\subsection{The Problem}

Modern engineering tools successfully manage artifacts.

Examples include:

\begin{itemize}
    \item Source code
    \item Design files
    \item Version history
    \item Documentation
    \item Simulation results
    \item Reports
\end{itemize}

However, these tools rarely preserve the reasoning that produced those artifacts.

Questions such as the following often remain unanswered:

\begin{itemize}
    \item Why was a particular architecture selected?
    \item Which alternatives were considered?
    \item What constraints influenced the decision?
    \item Why was an existing solution rejected?
    \item How did the design evolve over time?
\end{itemize}

As projects grow, this missing context becomes increasingly difficult to reconstruct.

\vspace{0.5cm}

\subsection{Our Beliefs}

We believe that:

\begin{enumerate}
    \item Engineering knowledge extends beyond source code.
    \item Design decisions are first-class engineering artifacts.
    \item Engineering context continuously evolves throughout the project lifecycle.
    \item Preserving reasoning is as important as preserving implementation.
    \item Documentation should describe understanding rather than merely recording activities.
    \item Engineering knowledge should remain accessible without relying on human memory.
    \item Knowledge captured once should be reusable throughout future projects.
\end{enumerate}

\vspace{0.5cm}

\subsection{Design Principles}

Every future design decision within this project should follow the principles below.

\begin{enumerate}
    \item Minimize interruption to the engineer's workflow.
    \item Preserve context rather than simply recording events.
    \item Prefer understanding over observation.
    \item Treat engineering decisions as persistent project knowledge.
    \item Allow engineering knowledge to evolve as designs evolve.
    \item Generate documentation from understanding rather than generating understanding from documentation.
    \item Maintain transparency regarding captured knowledge and generated outputs.
    \item Ensure that every generated artifact can be traced back to supporting engineering context.
\end{enumerate}

\vspace{0.5cm}

\subsection{Vision}

We envision engineering environments where project knowledge no longer disappears as engineers move between tasks, meetings, iterations, or team members.

Instead, engineering understanding continuously evolves alongside the project itself, enabling efficient collaboration, onboarding, design reviews, and long-term knowledge preservation.

\vspace{0.5cm}

\subsection{Success}

This project will be considered successful if it enables engineers to communicate project context, design evolution, and engineering reasoning with significantly less effort than is required using current engineering workflows.

Success is measured not by the amount of generated documentation, but by the quality, continuity, and accessibility of preserved engineering understanding.